\section*{Chapter \ref{ch:b1:induction} --- Electromagnetic Induction and Applications}

\begin{solution}{exo:b1:induction:1}
$\Phi = NBS\cos30^\circ = 100 \times 0.2 \times 10^{-3} \times 0.866 = \qty{1.7e-2}{Wb}$;
$\langle e\rangle = \Delta\Phi/\Delta t = \qty{1.7}{V}$ in \qty{10}{ms}, \qty{17}{V} in
\qty{1}{ms}. The current flows so as to recreate the vanishing flux: its
field is along the old $\vect B$.
\end{solution}

\begin{solution}{exo:b1:induction:2}
$e = B\ell v = 0.5 \times 0.2 \times 2 = \qty{0.20}{V}$; $i = \qty{2.0}{A}$; $F = B\ell i =
\qty{0.20}{N}$ braking, so the pusher applies \qty{0.20}{N}; $Fv = \qty{0.40}{W}
= Ri^2$.
\end{solution}

\begin{solution}{exo:b1:induction:3}
$L = \mu_0n^2S\ell = 4\pi \times 10^{-7} \times 10^6 \times 4 \times 10^{-4} \times 0.2 =
\qty{0.10}{mH}$; $W = \tfrac12LI^2 = \qty{1.3}{mJ}$; $\tau = L/R = \qty{0.1}{ms}$.
$B = \mu_0nI = \qty{6.3}{mT}$, $B^2/2\mu_0 = \qty{15.7}{J/m^3}$ over
$S\ell = \qty{8e-5}{m^3}$: \qty{1.3}{mJ}.
\end{solution}

\begin{solution}{exo:b1:induction:4}
$N_2/N_1 = 12/230 = 0.052$; $N_1 = 50/0.052 = 960$ turns; $i_1 = 2 \times 0.052 =
\qty{0.10}{A}$. A steady current makes a steady flux: no $\dd\varphi/\dd t$, no
secondary \omterm{def:b1:dc-circuits:voltage}{voltage} --- and the primary, a short, would burn.
\end{solution}

\begin{solution}{exo:b1:induction:5}
(a) The flux (along the approach) grows: the current makes a field
opposing it --- its own north pole faces the magnet's north --- and the
ring is pushed away. (b) Current whose field opposes the growth; the
ring, a dipole anti-aligned with $\vect B$, is pushed toward weaker
field. (c) As it falls the flux grows: same current as (a), the ring is
braked. (d) Each section of tube ahead of the magnet sees growing flux,
behind it decreasing; both induced currents brake the magnet; at the
speed where the \omterm{thm:b1:magnetostatics:laplace}{Laplace force} balances the weight, the magnet falls
steadily and its lost \omterm{def:b1:work-and-energy:conservative}{potential energy} becomes Joule heat in the tube.
\end{solution}

\begin{solution}{exo:b1:induction:6}
$E_{\max} = NBS\omega = 200 \times 2 \times 10^{-3} \times 0.1 \times 314 = \qty{12.6}{V}$, rms
\qty{8.9}{V}; $I_{\max} = \qty{1.26}{A}$, $\Gamma_{\max} = NBSI_{\max} = \qty{0.050}{N.m}$
(when $\sin\omega t = 1$); $\langle P\rangle = E_{\max}^2/2R = \qty{7.9}{W}$ $=
\langle\Gamma\omega\rangle = \tfrac12 \times 0.05 \times 314$; when $e = 0$ the current is
zero and the \omterm{def:b1:angular-momentum:moment}{torque} vanishes --- the \omterm{def:b1:angular-momentum:moment}{torque} pulses at \qty{100}{Hz}.
\end{solution}

\begin{solution}{exo:b1:induction:7}
$e = Bav$, $i = Bav/R$, $F = B^2a^2v/R = 10\,v$ (N). $m\,\dd v/\dd t = -B^2a^2v/R$:
$v = v_0\eu^{-t/\tau}$, $\tau = mR/B^2a^2 = \qty{0.10}{s}$; distance $v_0\tau =
\qty{10}{cm}$ (if the field region is long enough). The \omterm{thm:b1:work-and-energy:ke}{kinetic energy},
\qty{0.5}{J}, is dissipated as $Ri^2$ in the loop.
\end{solution}

\begin{solution}{exo:b1:induction:8}
Flux of the solenoid's field $\mu_0n_1i_1$ through the $N_2$ turns: $M =
\mu_0n_1N_2S_2 = 4\pi \times 10^{-7} \times 2000 \times 50 \times 2 \times 10^{-4} = \qty{25}{\micro H}$.
Ramp: $e = M\,\dd i_1/\dd t = \qty{2.5}{mV}$; at \qty{10}{kHz}, \qty{1}{A}: $e_{\max} =
M\omega I = 25 \times 10^{-6} \times 6.3 \times 10^4 = \qty{1.6}{V}$. $M$ is symmetric:
the same number gives the flux the small coil sends through the
solenoid --- far harder to compute directly.
\end{solution}

\begin{solution}{exo:b1:induction:9}
$W = B^2V/2\mu_0 = 2.25/(2.51 \times 10^{-6}) = \qty{0.90}{MJ}$; $L = 2W/I^2 =
\qty{7.2}{H}$; $h = W/mg = \qty{92}{m}$; helium: $9 \times 10^5/2600 = \qty{350}{L}$.
Earth's field: $(5 \times 10^{-5})^2/2\mu_0 = \qty{1e-3}{J/m^3}$, \qty{50}{mJ} in
the room.
\end{solution}

\begin{solution}{exo:b1:induction:10}
(a) $E - B\ell v = Ri$ (the back-emf opposes the source), $m\,\dd v/\dd t =
B\ell i$: $m\,\dd v/\dd t = B\ell(E - B\ell v)/R$. (b) $v_\infty = E/B\ell$, $\tau =
mR/B^2\ell^2$. (c) Multiply the electrical equation by $i$: $Ei = Ri^2 +
B\ell vi$, and $B\ell vi = Fv$ is the mechanical power; \omterm{def:b1:heat-engines:efficiency}{efficiency} $Fv/Ei =
B\ell v/E = v/v_\infty$ --- $100\%$ only at no load, where nothing is
delivered. (d) $v_\infty = 2/(0.5 \times 0.2) = \qty{20}{m/s}$, $\tau = 0.05 \times
0.1/0.01 = \qty{0.5}{s}$. (e) $F = B\ell I = 1 \times 0.1 \times 10^6 = \qty{1e5}{N}$
(taking $\ell = \qty{10}{cm}$), $a = \qty{1e7}{m/s^2}$, $v = \sqrt{2a \times 5} =
\qty{1e4}{m/s}$ --- ten kilometres per second, the promise and the
problem of rail guns (the rails erode).
\end{solution}

\begin{solution}{exo:b1:induction:11}
(a) At radius $r$ the metal moves at $v = \omega r$; $\vect v\wedge\vect B$ is
radial, of magnitude $\omega rB$: $e = \int_0^R\omega Br\,\dd r = \tfrac12B\omega R^2$.
(b) $\omega = 314$: $e = 0.5 \times 1 \times 314 \times 0.01 = \qty{1.6}{V}$. (c) One
``turn'' only, so volts at most; into \qty{1}{m\ohm}: \qty{1.6}{kA}, \omterm{def:b1:angular-momentum:moment}{torque}
$= P/\omega = ei/\omega = 2500/314 = \qty{8}{N.m}$. (d) The circuit brush--radius--
rim--brush sweeps out area $\tfrac12R^2\omega\,\dd t$ per $\dd t$: $\dd\Phi/\dd t =
\tfrac12BR^2\omega$ --- the flux through the moving, deforming circuit.
\end{solution}

\begin{solution}{exo:b1:induction:12}
(a) $e_2 = -M\,\dd i_1/\dd t$, peak $M\omega I_1 = 10^{-6} \times 1.57 \times 10^5 \times 30 =
\qty{4.7}{V}$. (b) $I_2 = 4.7/0.01 = \qty{470}{A}$ peak; $\langle P\rangle = R_pI_2^2/2
= \qty{1.1}{kW}$. (c) The emf is $\propto\omega$: high frequency gives a useful
emf with a modest $M$; the glass is an insulator, nothing flows in it
--- the heat is born inside the metal. (d) The pan current induces in
the coil $M\omega I_2 = 10^{-6} \times 1.57 \times 10^5 \times 470 = \qty{74}{V}$ peak,
in phase with $i_1$ (the pan current lags $e_2$ by nothing, being
resistive, and $e_2$ lags $i_1$ by \qty{90}{\degree}\ldots\ so the
reflected emf is $-M\,\dd i_2/\dd t$, in phase opposition to $i_1$'s
\emph{driving} \omterm{def:b1:dc-circuits:voltage}{voltage}): the generator supplies $\tfrac12 \times 74 \times 30 =
\qty{1.1}{kW}$ more --- exactly what the pan dissipates.
\end{solution}

\begin{solution}{pb:b1:induction:1}
\textbf{1.} $\Phi = NBS\cos\omega t$, $e = NBS\omega\sin\omega t$; $NBS = 300 \times
0.3 \times 3 \times 10^{-4} = \qty{0.027}{Wb}$; $\omega = v/r = 5/0.01 = \qty{500}{rad/s}$:
\omterm{def:b1:wave-propagation:signal}{amplitude} \qty{13.5}{V}, frequency $\omega/2\pi = \qty{80}{Hz}$.

\textbf{2.} $I_{\text{rms}} = 13.5/\sqrt2/16 = \qty{0.60}{A}$; lamp $RI^2 = \qty{4.3}{W}$
under \qty{7.2}{V}: a \qty{6}{V}, \qty{3}{W} lamp is over-driven and will not
last.

\textbf{3.} Mechanical $=$ electrical: $(R + r_c)I_{\text{rms}}^2 = 16 \times 0.36 =
\qty{5.7}{W}$; $F = P/v = 5.7/5 = \qty{1.1}{N}$ --- half the rolling friction;
the cyclist feels it.

\textbf{4.} The induced current's Laplace couple opposes the rotation
(Lenz); the lamp's energy comes from the cyclist's legs, \qty{5.7}{W} of it.

\textbf{5.} $\underline Z = R + r_c + jL\omega$; $I = NBS\omega/\sqrt{(R + r_c)^2 + L^2\omega^2}$;
for $L\omega \gg R + r_c$, $I \to NBS/L = 0.027/0.02 = \qty{1.35}{A}$. $70\%$: $L\omega
= 0.7(R + r_c)/\sqrt{1 - 0.49} = 15.7$, $\omega = \qty{784}{rad/s}$, $v = \qty{7.8}{m/s}
= \qty{28}{km/h}$.

\textbf{6.} Without $L$: $e_{\max} = 3 \times 13.5 = \qty{40.5}{V}$, $I_{\text{rms}} =
\qty{1.8}{A}$, \qty{39}{W} in the lamp --- it blows. With $L$: $I = 40.5/\sqrt{256 +
900} = \qty{1.2}{A}$ peak, \qty{8.5}{W}: hot, but ten times less.

\textbf{7.} \qty{9}{km/h}: $e_{\max} = 6.75$, $|Z| = \sqrt{256 + 25} = 16.8$, $I = 0.40$,
$P = \tfrac12 \times 0.16 \times 12 = \qty{1.0}{W}$: dim. \qty{36}{km/h}: $e_{\max} = 27$,
$|Z| = 25.6$, $I = 1.05$, $P = \qty{6.6}{W}$. Four times the speed, $2.6$ times
the current: the inductance regulates.

\textbf{8.} Only the \omterm{def:b1:non-inertial-frames:frames}{relative motion} matters: the flux through the fixed
coil varies identically, same emf. No sliding contacts (brushes) are
needed to bring the current out of a rotating coil.

\textbf{9.} No motion, no flux change, no light --- hence the \omterm{def:b1:potential-capacitors:capacitance}{capacitor} or
small battery (``standlight'') charged while riding.

\textbf{10.} Each element of wire is perpendicular to the radial $\vect B$:
$F = B\ell i$, axial. Moving at $v$, each element sees $\vect v\wedge\vect B$
along the wire: $e = -B\ell v$, opposing the current that produces the
motion (\omterm{thm:b1:dc-circuits:kirchhoff}{receiver convention}: a \omterm{def:b1:dc-circuits:voltage}{voltage} $B\ell v$).

\textbf{11.} $u = Ri + L\,\dd i/\dd t + B\ell v$; $m\,\dd v/\dd t = -kx - hv + B\ell i$.

\textbf{12.} $\underline i = (\underline u - B\ell\underline v)/(R + jL\omega)$; $\underline v\,(jm\omega + h +
k/j\omega) = B\ell\,\underline i$; substitute and solve for $\underline v$.

\textbf{13.} $i \approx (u - B\ell v)/R$ gives $m\dot v + (h + B^2\ell^2/R)v + kx =
B\ell u/R$: extra damping $25/8 = \qty{3.1}{N.s/m}$. $f_0 = \tfrac1{2\pi}\sqrt{k/m} =
\qty{50}{Hz}$; $Q = \sqrt{km}/h_{\text{tot}}$: $3.2$ open, $0.77$ with the
amplifier.

\textbf{14.} $F = \qty{5.0}{N}$, $a = F/m = \qty{500}{m/s^2}$, $x = a/\omega^2 =
\qty{13}{\micro m}$, $v = a/\omega = \qty{0.080}{m/s}$; $B\ell v = \qty{0.40}{V}$ against
$Ri = \qty{8}{V}$: $5\%$.

\textbf{15.} With the current imposed, $v = F/h = \qty{5}{m/s}$ (an
\omterm{def:b1:wave-propagation:signal}{amplitude} of \qty{16}{mm} --- a big excursion for \qty{1}{A}); $B\ell v =
\qty{25}{V}$, three times $Ri$: the amplifier must supply \qty{33}{V} to push
\qty{1}{A} at resonance --- the \omterm{def:b1:sinusoidal-impedance:impedance}{impedance} peak. A voltage-source amplifier
would instead let the electrical damping hold the cone.

\textbf{16.} $\tfrac12RI^2 = \qty{4.0}{W}$; $\tfrac12hv^2 = \tfrac12 \times 0.0064 =
\qty{3.2}{mW}$: \omterm{def:b1:heat-engines:efficiency}{efficiency} $0.08\%$ --- loudspeakers are heaters that
whisper.

\textbf{17.} Above resonance $a = F/m = B\ell i/m$ is independent of
frequency for a given current, so the pressure, $\propto a$, is flat:
the mass-controlled band is the working band.

\textbf{18.} Radial field: the force is axial and the same at every
position of the coil in the gap. A tweeter must stay mass-controlled
to \qty{20}{kHz} with a tiny excursion: light coil and cone; and a
small cone radiates high frequencies in all directions.

\textbf{19.} Low frequency: $v$ small, $|Z| \approx R = \qty{8}{\ohm}$; at
\qty{50}{Hz}: $u = Ri + B\ell v$ with $v = B\ell i/h$: $Z = R + B^2\ell^2/h = 8 + 25 =
\qty{33}{\ohm}$; at \qty{10}{kHz}: $|R + jL\omega| = |8 + 31j| = \qty{32}{\ohm}$. A
peak at resonance, a flat \qty{8}{\ohm} floor, a rise with $L\omega$.

\textbf{20.} $e = B\ell v = 5 \times 10^{-3} = \qty{5}{mV}$.

\textbf{21.} $i = 5 \times 10^{-3}/608 = \qty{8.2}{\micro A}$; $F = B\ell i = \qty{41}{\micro N}$,
opposing $v$: $F = [B^2\ell^2/(R + R_L)]\,v$, a damping of \qty{0.041}{N.s/m};
the power $Fv = \qty{4e-8}{W}$ is taken from the sound and ends in the
two \omterm{def:b1:dc-circuits:resistor}{resistors}.

\textbf{22.} Former: $B^2\ell_f^2/R_f = 1 \times 0.01/10^{-3} = \qty{10}{N.s/m}$, ten
times $h$: it would deaden the cone and waste power as heat; a slit
breaks the closed turn.

\textbf{23.} Open: only $h$ damps, $Q = 3.2$, it rings. Shorted: the
back-emf drives a current through $R$, damping $B^2\ell^2/R$ is added,
$Q = 0.77$: it stops at once. (Try it.)

\textbf{24.} Over a period, stored energies return to their values:
$\langle ui\rangle = \langle Ri^2\rangle + \langle hv^2\rangle$ --- electrical input $=$
Joule heat in the coil $+$ mechanical losses (sound radiated and
suspension friction); the coupling term $B\ell iv$ appears in both
equations with opposite signs and cancels.

\textbf{25.} Dynamo: Faraday on a rotating coil --- \qty{13.5}{V} at
\qty{18}{km/h}, current capped at \qty{1.35}{A} by its own inductance.
Loudspeaker: \omterm{thm:b1:magnetostatics:laplace}{Laplace force} and back-emf --- \qty{3.1}{N.s/m} of electrical
damping, $0.08\%$ \omterm{def:b1:heat-engines:efficiency}{efficiency}. Microphone: the same coil backward ---
\qty{5}{mV} per millimetre per second.
\end{solution}
